\section{記法と量子力学の公理}
前節の破綻は、少なくとも次の三点を古典力学とは別の仕方で定め直すことを要求している。すなわち、状態を何で表すか、物理量を何で表すか、そして時間発展をどう与えるか、である。歴史的にはこれらは個別の現象を通じて段階的に見出されたが、本ノートでは以後の議論を見通しよくするため、まず公理的枠組みとしてまとめて提示する。

\subsection{状態（公理I）}
系の物理的状態は、複素ヒルベルト空間の元（状態ベクトル）$\psi$ によって表される。ヒルベルト空間とは、内積 $\langle \psi,\varphi\rangle$ が定義された複素ベクトル空間であって、その内積から定まるノルムに関して完備なものである。量子論で複素数を用いるのは、位相差が干渉を支配するからである。$\psi$ が複素係数倍（全体位相）の違いを除いて同一なら、物理的に同一の状態である。ノルムが1になるような係数を選ぶことを「正規化」といい、
\[
\langle \psi,\psi\rangle =1
\]
を満たすようにする。

\subsection{観測量（公理II）}
物理量はヒルベルト空間上の自己共役演算子 $\hat{A}$ によって表される。本ノートでは、演算子であることを明示するためにハット（$\hat{\quad}$）を付す。自己共役であることは、測定値が実数になることに対応している。

\subsection{測定と確率（公理III）}
状態 $\psi$（ただし正規化されているとする）において物理量 $\hat{A}$ を測定したとき、得られる値は $\hat{A}$ の固有値のいずれかであり、その確率は状態ベクトルと固有状態の内積の絶対値の2乗（ボルンの規則）によって与えられる。

\subsection{時間発展（公理IV）}
時間発展は線形であり、かつ確率（内積）を保存するユニタリ変換として与えられる。これは、全確率
\[
\langle \psi(t),\psi(t)\rangle
\]
が時間に依らないことを意味する。具体的なヒルベルト空間の選び方は、各理論を導入する段階でその都度与える。
